\documentclass[aspectratio=169]{beamer}

\usetheme{Madrid}
\usecolortheme{beaver}

\usepackage{graphicx}
\usepackage{booktabs}
\usepackage{qrcode}
\usepackage{hyperref}
\usepackage{fontawesome5}

\title[Event Countdown Timer]{Event Countdown Timer}
\subtitle{A real-time event countdown with desktop notifications}
\author{Fedor Kozhevnikov}
\institute[Uni]{Innopolis University \\ \texttt{f.kozhevnikov@innopolis.university} \\ Group: CSE-04}
\date{\today}

\begin{document}

% Slide 1: Title
\begin{frame}
    \titlepage
\end{frame}

% Slide 2: Context
\begin{frame}{Context}
    \begin{block}{End-User}
        Students, professionals, and anyone who tracks deadlines and events
    \end{block}

    \begin{block}{Problem}
        People miss important events because they forget to check calendars or don't get timely reminders
    \end{block}

    \begin{block}{Product One-Liner}
        \centering
        \textit{``A web page where you add events, watch live countdowns, and get desktop notifications when they arrive''}
    \end{block}
\end{frame}

% Slide 3: Implementation
\begin{frame}{Implementation}
    \small
    \setlength{\itemsep}{0pt}
    \setlength{\parskip}{0pt}

    \textbf{Tech Stack:} Python + FastAPI, SQLite, APScheduler, Vanilla HTML/CSS/JS

    \vspace{0.1cm}

    \begin{block}{Version 1 -- Core MVP}
        \begin{itemize}
            \item Add/delete events with title and target date/time
            \item Real-time countdown (days, hours, minutes, seconds)
            \item Persistent storage via SQLite
            \item Clean, responsive UI
        \end{itemize}
    \end{block}

    \begin{block}{Version 2 -- Notifications \& Polish}
        \begin{itemize}
            \item Desktop push notifications on event start
            \item Background scheduler for reliable delivery
            \item Visual state changes (color shifts) as events approach
            \item Server-client time synchronization
            \item User session isolation (private event views)
        \end{itemize}
    \end{block}
\end{frame}

% Slide 4: Demo (placeholder)
\begin{frame}{Demo}
    \begin{center}
        \vspace{1cm}
        {\Huge \textbf{Demo}}

        \vspace{0.5cm}
        {\Large Pre-recorded video demonstration of Version 2}

        \vspace{1cm}
        \textit{[Insert your 2-minute video with voice-over here]}
    \end{center}
\end{frame}

% Slide 5: Links
\begin{frame}{Links}
    \begin{center}
        % First Column: GitHub Repository
        \begin{minipage}[t]{0.45\textwidth}
            \centering
            \textbf{GitHub Repository}

            \vspace{0.2cm}
            \href{https://github.com/Fedos113/se-toolkit-hackathon}{\texttt{github.com/Fedos113/se-toolkit-hackathon}}

            \vspace{0.3cm}
            \qrcode[height=3cm]{https://github.com/Fedos113/se-toolkit-hackathon}
        \end{minipage}
        \hfill
        % Second Column: Deployed Product
        \begin{minipage}[t]{0.45\textwidth}
            \centering
            \textbf{Deployed Product (Latest Version)}

            \vspace{0.2cm}
            \href{http://10.93.25.224:8000}{\texttt{http://10.93.25.224:8000}}

            \vspace{0.3cm}
            % Note: Removed duplicate 'http://' in the QR code data which was present in your original snippet
            \qrcode[height=3cm]{http://10.93.25.224:8000}
        \end{minipage}
    \end{center}
\end{frame}

\end{document}
